\SourceLeaf{162}
\SixthArticle{30}
Aus Formel 38 folgt
\[
\frac{\Pi(k,z+\omega)}{\Pi(k,z)}
=\frac{z+1}{z+1+\omega}\cdot\frac{z+2}{z+2+\omega}\cdot\frac{z+3}{z+3+\omega}
\mathinner{.\,.\,.\,.}\frac{z+k}{z+k+\omega}\cdot k^{\omega}
\]
Nimmt man die Logarithmen und entwickelt sie in unendliche Reihen, so ergibt sich
\[
\begin{aligned}
{}\text{[60]}\quad&\log\Pi(k,z+\omega)=\log\Pi(k,z)\\
&+\omega\left(\log k-\frac{1}{z+1}-\frac{1}{z+2}-\frac{1}{z+3}-\text{usw.}-\frac{1}{z+k}\right)\\
&+\tfrac12\omega\omega\left(\frac{1}{(z+1)^2}+\frac{1}{(z+2)^2}+\frac{1}{(z+3)^2}+\text{usw.}+\frac{1}{(z+k)^2}\right)\\
&-\tfrac13\omega^3\left(\frac{1}{(z+1)^3}+\frac{1}{(z+2)^3}+\frac{1}{(z+3)^3}+\text{usw.}+\frac{1}{(z+k)^3}\right)\\
&+\text{usw. ins Unendliche}
\end{aligned}
\]
Die hier mit $\omega$ multiplizierte Reihe, die, wenn man es vorzieht, auch so dargestellt werden kann:
\[
\begin{aligned}
&-\frac{1}{z+1}+\log 2-\frac{1}{z+2}+\log\tfrac32-\frac{1}{z+3}\\[4pt]
&+\log\tfrac43-\frac{1}{z+4}+\log\tfrac54-\text{usw.}
+\log\frac{k}{k-1}-\frac{1}{z+k}
\end{aligned}
\]
besteht aus einer endlichen Anzahl von Gliedern; wächst aber $k$ ins Unendliche, so konvergiert sie gegen einen bestimmten Grenzwert, der uns eine neue Art transzendenter Funktionen liefert, die künftig mit $\Psi z$ bezeichnet werden soll.

Bezeichnen wir ferner die Summen der folgenden, ins \textit{Unendliche} fortgesetzten Reihen
\[
\begin{array}{l}
\dfrac{1}{(z+1)^2}+\dfrac{1}{(z+2)^2}+\dfrac{1}{(z+3)^2}+\text{usw.}\\[10pt]
\dfrac{1}{(z+1)^3}+\dfrac{1}{(z+2)^3}+\dfrac{1}{(z+3)^3}+\text{usw.}\\[10pt]
\dfrac{1}{(z+1)^4}+\dfrac{1}{(z+2)^4}+\dfrac{1}{(z+3)^4}+\text{usw.}\\[10pt]
\text{usw.}
\end{array}
\]
der Reihe nach mit $P,Q,R$ usw. (für die es weniger notwendig scheint, Funktionszeichen einzuführen), so haben wir
\[
\begin{aligned}
{}\text{[61]}\quad&\log\Pi(z+\omega)=\log\Pi z+\omega\Psi z\\[4pt]
&+\tfrac12\omega\omega P-\tfrac13\omega^3Q+\tfrac14\omega^4R-\text{usw.}
\end{aligned}
\]
Offenbar ist die Funktion $\Psi z$ die erste Ableitung der Funktion $\log\Pi z$, und daher ist
\[
{}\text{[62]}\qquad \frac{\mathrm{d}\Pi z}{\mathrm{d}z}=\Pi z\cdot\Psi z
\]
Ebenso ist $P=\dfrac{\mathrm{d}\Psi z}{\mathrm{d}z}$, $Q=-\dfrac{\mathrm{dd}\Psi z}{2\mathrm{d}z^2}$, $R=+\dfrac{\mathrm{d}^3\Psi z}{2.3\mathrm{d}z^3}$ usw.
\SourceLeaf{163}
\SixthArticle{31}
Die Funktion $\Psi z$ ist fast ebenso bemerkenswert wie die Funktion $\Pi z$; deshalb wollen wir hier die wichtigeren sie betreffenden Beziehungen zusammenstellen. Durch Differentiation der Gleichung 44 ergibt sich
\[
\Psi(z+1)=\Psi z+\frac{1}{z+1}\tag*{[63]}
\]
woraus folgt
\[
\Psi(z+n)=\Psi z+\frac{1}{z+1}+\frac{1}{z+2}+\frac{1}{z+3}+\text{usw.}+\frac{1}{z+n}\tag*{[64]}
\]
Mit ihrer Hilfe kann man von kleineren Werten von $z$ zu größeren fortschreiten oder von größeren zu kleineren zurückgehen. Für größere positive Werte von $z$ lassen sich die Zahlenwerte der Funktion recht bequem durch die folgenden Formeln berechnen, die aus der Differentiation der Gleichungen 58 und 59 hervorgehen; dabei ist jedoch dasselbe zu beachten, worauf wir in Artikel 29 bei den Formeln 58 und 59 hingewiesen haben:
\[
\Psi z=\log z+\frac{1}{2z}-\frac{\mathfrak{A}}{2zz}+\frac{\mathfrak{B}}{4z^4}-\frac{\mathfrak{C}}{6z^6}+\text{usw.}\tag*{[65]}
\]
\[
\Psi\left(z-\tfrac12\right)=\log z+\frac{\mathfrak{A}}{2\cdot2zz}-\frac{7\mathfrak{B}}{4\cdot8z^4}+\frac{31\mathfrak{C}}{6\cdot32z^6}-\text{usw.}\tag*{[66]}
\]
So haben wir für $z=10$ berechnet:
\[
\Psi z=2{,}3517525890\ 6672110764\ 743
\]
woraus wir zurückgehen zu
\[
\Psi0=-0{,}5772156649\ 0153286060\ 653\HistoricalNote{18}
\]
Für einen positiven ganzzahligen Wert von $z$ gilt allgemein
\[
\Psi z=\Psi0+1+\tfrac12+\tfrac13+\text{usw.}+\frac1z\tag*{[67]}
\]
Für einen negativen ganzzahligen Wert wird $\Psi z$ dagegen offenbar zu einer unendlich großen Größe.
\SourceLeaf{164}
\SixthArticle{32}
Formel 55 liefert uns $\log\Pi(-z)+\log\Pi(z-1)=\log\pi-\log\sin z\pi$, woraus durch Differentiation folgt:
\[
\Psi(-z)-\Psi(z-1)=\pi\operatorname{cotang}z\pi\tag*{[68]}
\]
Und da nach der Definition der Funktion $\Psi$ allgemein
\[
\Psi x-\Psi y=-\frac1{x+1}+\frac1{y+1}-\frac1{x+2}+\frac1{y+2}-\frac1{x+3}+\text{usw.}\tag*{[69]}
\]
gilt, ergibt sich die bekannte Reihe
\[
\pi\operatorname{cotang}z\pi=\frac1z-\frac1{1-z}+\frac1{1+z}-\frac1{2-z}+\frac1{2+z}-\frac1{3-z}+\text{usw.}
\]
Ebenso ergibt sich durch Differentiation der Formel 57
\[
\begin{aligned}
&\Psi z+\Psi\left(z-\frac1n\right)+\Psi\left(z-\frac2n\right)+\text{usw.}+\Psi\left(z-\frac{n-1}{n}\right)\\[4pt]
&=n\Psi nz-n\log n
\end{aligned}
\tag*{[70]}\]
und daher, wenn $z=0$ gesetzt wird,
\[
\begin{aligned}
&\Psi\left(-\frac1n\right)+\Psi\left(-\frac2n\right)+\Psi\left(-\frac3n\right)+\text{usw.}+\Psi\left(-\frac{n-1}{n}\right)\\[4pt]
&=(n-1)\Psi0-n\log n
\end{aligned}
\tag*{[71]}\]
So gilt beispielsweise
\[
\begin{aligned}
&\Psi\left(-\tfrac12\right)=\Psi0-2\log2=-1{,}9635100260\ 2142347944\ 099,\\[4pt]
&\quad\text{woraus ferner}
\end{aligned}
\]
\[
\Psi\tfrac12=+0{,}0364899739\ 7857652055\ 901.
\]
